\section{Finite rational decision tables}\label{sec:finite-tables}

This section converts the companion existence theorem into a terminating
finite computation. The output is a decision table for every prescribed
horizon, with dyadic probabilities and a bounded total rounding loss.
No running-time bound on this finite computation is needed yet.

Throughout, fix the input metric on $X=\{1,\ldots,n\}$ and put
\begin{equation}\label{eq:metric-parameters}
 D=\max_{x,y\in X}d(x,y),\qquad
 \delta=\min_{x\ne y}d(x,y),\qquad u=(1,\ldots,k).
\end{equation}
Thus $0<\delta\le D$, and the auxiliary initial tuple $u$ has distinct
entries, even when the actual initial tuple $s$ does not.
For a request word $w$ and initial labeled tuple $v\in X^k$, write
$\operatorname{OPT}_v(w)$ for the offline minimum. If $|w|=h$, this
minimum is attained among the $k^h$ sequences of serving labels.
Allowing an offline comparator to move other servers between requests
cannot improve this minimum: choose a serving label at each request,
then move each label directly along its assigned subsequence of requests.
The triangle inequality makes the resulting service no more expensive.
We will use this shortcutting observation again when deleting requests.

\subsection{The existence input and strong laziness}

The only theorem imported into the construction is the following
specialization of the companion's main result. We use its distinct-start
clause, which has zero additive loss.

\begin{theorem}[Existence from distinct initial positions,
 {\cite[Theorem~1.1]{SquaredLog2026}}]\label{thm:existence}
There is an absolute constant $C_0<\infty$ such that, for every integer
$k\ge2$, every finite metric space with at least $k+1$ points, and every
distinct initial labeled tuple $u$, there exists a randomized online
policy $A$ satisfying
\begin{equation}\label{eq:external-existence}
 \E\bigl[\operatorname{cost}_{A,u}(w)\bigr]
 \le C_0\bigl(\log(k+1)\bigr)^2\operatorname{OPT}_u(w)
\end{equation}
for every fixed finite request word $w$ chosen independently of its
random choices. The same policy serves all words and all horizons. At
each request it chooses one labeled server and moves that server to the
request, leaving the others fixed.
\end{theorem}

The cited theorem applies to arbitrary metric spaces and allows repeated
initial positions with a finite additive term; its distinct-position
conclusion is exactly the zero-additive assertion stated here. In our
application $2\le k<n$ and $u=(1,\ldots,k)$, so all its hypotheses hold.
Only policy existence is used: neither its probabilities nor a procedure
for computing them is supplied to the construction below.

Call a policy \emph{strongly lazy} if, whenever the request is already
occupied, it chooses a server at that request and pays zero. On a miss
it moves one server to the request. From a distinct tuple, such a policy
keeps the positions distinct.

The following is the standard lazification principle of Manasse, McGeoch,
and Sleator~\cite[Lemma~1]{MMS1990}. We include a matching-potential proof
to make the labeled, pathwise comparison explicit.

\begin{lemma}\label{lem:lazy}
From a distinct initial tuple, any online policy has a strongly lazy
online follower whose total cost on every finite word is at most the
policy's cost, for every realization of the policy's randomness.
\end{lemma}

\begin{proof}
Simulate the original policy virtually. For follower tuple $c$ and
virtual tuple $v$, use the matching potential
\[
 \Phi(c,v)=\min_{\pi\in S_k}\sum_{j=1}^k d(c_j,v_{\pi(j)}).
\]
The minimum exists even if the virtual positions repeat. Initially the
potential is zero. Suppose the virtual policy has just served a request
$r$, at cost $\eta$, changing $v$ to $v'$. The triangle inequality,
applied to a previously minimizing matching, gives
$\Phi(c,v')\le\Phi(c,v)+\eta$.

If $c$ covers $r$, the follower makes a zero-cost choice there. Otherwise
choose a minimizing matching between $c$ and $v'$ and a virtual label
at $r$, and move its matched follower label to $r$. If this move costs
$m$ and gives tuple $c'$, its matching edge decreases from length $m$
to zero, while all other edges remain unchanged. Hence in either case
\[
 m+\Phi(c',v')\le\Phi(c,v)+\eta.
\]
Summing and using nonnegativity of the terminal potential proves cost
domination. The follower moves only to an uncovered point, so its
positions remain distinct. Fixed orders for labels and permutations
resolve ties, and all choices use only the current simulated state and
the revealed request. Thus the construction is causal.
\end{proof}

\subsection{A finite flow system}

Finite-horizon linear descriptions of randomized $k$-server strategies
appear in M\"omke~\cite[Section~4]{Moemke2013}. We use configuration
marginals and transition flows, with a single marginal shared by all
continuations of each request prefix. The proposition also records the
finite-bit sampling guarantee needed later.

\begin{proposition}\label{prop:finite-table}
For every positive integer $H$, a terminating deterministic procedure
on rational input constructs a rational number $a$ and a finite table
defining a strongly lazy policy $\mathcal T_H$ from $u$ such that
\begin{equation}\label{eq:finite-table-bound}
 0\le a\le C_0\bigl(\log(k+1)\bigr)^2,\qquad
 \E\bigl[\operatorname{cost}_{\mathcal T_H,u}(w)\bigr]
 \le a\operatorname{OPT}_u(w)+D
 \quad (|w|\le H).
\end{equation}
Every table decision uses exactly $b$ independent unbiased bits, where
$b$ is the least nonnegative integer with $2^b\ge kH^2$. The procedure
requires neither the value of $C_0$ nor an oracle for the policy in
Theorem~\ref{thm:existence}.
\end{proposition}

\begin{proof}
Let $\mathcal C$ be the set of ordered $k$-tuples of distinct points.
For $c\in\mathcal C$ and $r\in X$, define the allowed labels by
\[
 J(c,r)=
 \begin{cases}
  \{j:c_j=r\},&r\in\{c_1,\ldots,c_k\},\\
  \{1,\ldots,k\},&r\notin\{c_1,\ldots,c_k\}.
 \end{cases}
\]
For any tuple $c\in X^k$, let $T(c,r,j)$ replace its $j$th entry by
$r$. If $c\in\mathcal C$ and $j\in J(c,r)$, the resulting tuple lies
in $\mathcal C$. All finite index sets below are ordered by length and
then lexicographically, so all otherwise arbitrary choices can be
deterministic.

For every $w\in X^{\le H}$ and $c\in\mathcal C$, introduce a
variable $p_{w,c}$, intended to be the probability of configuration $c$
after the word $w$. For $|w|<H$, $r\in X$, $c\in\mathcal C$, and
$j\in J(c,r)$, introduce a variable $f_{wr,c,j}$. Along with the
coefficient variable $a$, these are the unknowns of the finite system.
The intended meaning of $f_{wr,c,j}$ is the joint probability that the
configuration before request $r$ is $c$ and the serving label is $j$.
Impose the following constraints:
\begin{align}
 &p_{w,c}\ge0,\qquad \sum_{c\in\mathcal C}p_{w,c}=1,
 \qquad p_{\varnothing,c}=\mathbf1_{\{c=u\}},
 && w\in X^{\le H},                                      \label{eq:table-masses}\\
 &f_{wr,c,j}\ge0,\qquad
 \sum_{j\in J(c,r)}f_{wr,c,j}=p_{w,c},
 && |w|<H,                                               \label{eq:table-outflow}\\
 &\sum_{\substack{c\in\mathcal C,\ j\in J(c,r)\\
                   T(c,r,j)=c'}} f_{wr,c,j}=p_{wr,c'},
 && |w|<H,\ c'\in\mathcal C.                             \label{eq:table-inflow}
\end{align}
Here and below the request $r$ ranges over $X$, and every displayed
flow variable has an allowed label. Finally impose $a\ge0$ and, for
each nonempty word $x=r_1\cdots r_h$ with $h\le H$, impose
\begin{equation}\label{eq:table-competitiveness}
 \sum_{i=1}^h\sum_{c\in\mathcal C}\sum_{j\in J(c,r_i)}
     f_{r_1\cdots r_i,c,j}\,d(c_j,r_i)
 \le a\operatorname{OPT}_u(x).
\end{equation}

All coefficients are computable rationals. To make the last assertion
explicit, let $\Gamma_w(c)$ be the minimum cost of serving $w$ from
$u$ and finishing at $c\in X^k$, with $+\infty$ marking an unreachable
tuple. Start with $\Gamma_{\varnothing}(u)=0$ and all other entries
$+\infty$, and use
\[
 \Gamma_{wr}(c')=
 \min_{\substack{c\in X^k,\ 1\le j\le k\\T(c,r,j)=c'}}
        \bigl(\Gamma_w(c)+d(c_j,r)\bigr),
 \qquad
 \operatorname{OPT}_u(w)=\min_{c\in X^k}\Gamma_w(c).
\]
These finite recurrences compute every required coefficient by rational
arithmetic and comparisons. No division by $\operatorname{OPT}_u(w)$
is used; words with zero optimum simply give a zero upper bound in
Equation~\eqref{eq:table-competitiveness}.

By Theorem~\ref{thm:existence} and Lemma~\ref{lem:lazy}, the system has
a real solution with $a=C_0(\log(k+1))^2$: take $p_{w,c}$ to be the
configuration probabilities of the strongly lazy policy after $w$,
and $f_{wr,c,j}$ the probabilities of configuration $c$ immediately
before $r$ and serving label $j$. Causality gives one distribution
$p_{w,\cdot}$ shared by every continuation of $w$.

\paragraph{Exact rational attainment.}
We now compute a rational solution having the smallest feasible value
of $a$. Replace equalities by pairs of weak inequalities and eliminate
all variables except $a$ by Fourier--Motzkin elimination
\cite{DantzigEaves1973}. We record the exact projection and rational
back-substitution because termination, rather than efficiency, is the
property required here. To eliminate a variable $z$, divide each inequality
containing it by
its nonzero rational coefficient to obtain finitely many bounds
\[
 z\ge L_i(y),\qquad z\le U_j(y),
\]
where $y$ lists the remaining variables and the bounds are rational
affine expressions. Keep inequalities independent of $z$ and add all
comparisons $L_i(y)\le U_j(y)$. This is an exact existential projection:
if both bound families are present, their intervals intersect precisely
when every such comparison holds; if either family is absent, the
remaining one-sided bounds always admit a value of $z$.

After finitely many eliminations the feasible $a$-values are given by
finitely many weak rational inequalities in one variable. This set is
nonempty and includes the constraint $a\ge0$, so its least value is
the maximum of its finitely many rational lower bounds. It is rational
and belongs to the feasible set. Compute this value, then reverse the
eliminations. At each reversal, evaluated bounds are rational: choose
the largest lower bound if one exists; otherwise choose the smallest
upper bound if one exists; otherwise choose zero. The projection
inequalities ensure that the chosen value also respects any upper
bounds. This recovers a rational solution of the full system.

Thus all steps terminate on finite rational data, and the computed
minimum satisfies $0\le a\le C_0(\log(k+1))^2$. The constant $C_0$
never enters the computation; the existence theorem only certifies
feasibility and the upper bound on the minimum.

\paragraph{Causal conditional rows.}
For each $(w,r,c)$ define a probability distribution on $J(c,r)$ by
\begin{equation}\label{eq:table-conditional}
 q_{w,r,c}(j)=\frac{f_{wr,c,j}}{p_{w,c}}
 \quad\text{if }p_{w,c}>0.
\end{equation}
If $p_{w,c}=0$, instead put all mass on the smallest allowed label.
Every row is now a rational probability distribution, including rows
at configurations not reached by the unrounded policy. Choose from
the row indexed by the revealed prefix, next request, and current
configuration. Along any fixed word, induction using
Equations~\eqref{eq:table-outflow} and~\eqref{eq:table-inflow} gives
exactly the marginals $p$ and $f$: at a positive-mass configuration,
$p_{w,c}q_{w,r,c}(j)=f_{wr,c,j}$, and at a zero-mass configuration all
outgoing flows vanish. Consequently Equation~\eqref{eq:table-competitiveness}
is the expected-cost bound for these causal decisions. The construction
need not retain any hidden memory of the source policy.

\paragraph{Dyadic rows and total rounding loss.}
In each row reserve its smallest allowed label $j_0$. For $j\ne j_0$
set
\[
 \widehat q_{w,r,c}(j)=2^{-b}\lfloor 2^bq_{w,r,c}(j)\rfloor,
 \qquad
 \widehat q_{w,r,c}(j_0)=1-\sum_{j\ne j_0}\widehat q_{w,r,c}(j).
\]
The new row is supported on allowed labels, and its total variation
distance from the old row is
\[
 \frac12\sum_{j\in J(c,r)}
        |\widehat q_{w,r,c}(j)-q_{w,r,c}(j)|
 =\sum_{j\ne j_0}\bigl(q_{w,r,c}(j)-\widehat q_{w,r,c}(j)\bigr)
 \le k2^{-b}.
\]
Store each row by its integer numerators with common denominator $2^b$.
For sampling, use $b$ unbiased bits to choose a uniform integer in
$\{0,\ldots,2^b-1\}$ and return the first label whose cumulative
numerator exceeds it. This uses exactly $b$ random bits.

Fix a word $w$ of length $h\le H$. Couple the rounded and unrounded
policies until their first different label choice: whenever their
states agree, match each possible label with probability equal to the
minimum of its two row probabilities, and couple the residual masses
arbitrarily. At any such step the conditional probability of a
discrepancy is at most $k2^{-b}$. Thus the probability of any discrepancy
on the word is at most $hk2^{-b}$. On its complement the two costs
agree; on the discrepancy event their upward difference is at most
$hD$, since each total cost lies in $[0,hD]$. Therefore
\[
 \E\bigl[\operatorname{cost}_{\mathcal T_H,u}(w)\bigr]
 \le a\operatorname{OPT}_u(w)+k h^2 2^{-b}D
 \le a\operatorname{OPT}_u(w)+D.
\]
A rounded trajectory can reach a configuration with original
$p_{w,c}=0$. Such a row has already been defined legally, and reaching
it can only require attention after a discrepancy, where the uniform
bound $hD$ suffices. Hence these configurations cause no exception to
the proof or to the sampling rule. This proves the proposition.
\end{proof}

Proposition~\ref{prop:finite-table} gives a finite, fully specified
constructor. Its running time and table size may be enormous. The
remaining argument first turns one fixed table into a routine for an
unbounded suffix, and then schedules the constructor within the allowed
bit budget.
